% Chapter 2.10: Automata (2-OS) - LaTeX Beamer Slides
% Structured Computer Architecture: An Order-Based Approach
% Authors: John Glossner and Gheorghe M. Ştefan

\documentclass[aspectratio=169,11pt]{beamer}

% Theme and color scheme (simple, print-friendly)
\usetheme{default}
\usecolortheme{default}

% Simple color scheme for printing
\definecolor{titlecolor}{RGB}{0,0,0}
\definecolor{accentcolor}{RGB}{0,0,139}
\definecolor{textcolor}{RGB}{0,0,0}

% Set colors
\setbeamercolor{title}{fg=titlecolor}
\setbeamercolor{frametitle}{fg=titlecolor}
\setbeamercolor{structure}{fg=accentcolor}
\setbeamercolor{normal text}{fg=textcolor}

% Remove navigation symbols
\setbeamertemplate{navigation symbols}{}

% Simple footline
\setbeamertemplate{footline}[frame number]

% Packages
\usepackage{amsmath}
\usepackage{amsfonts}
\usepackage{amssymb}
\usepackage{graphicx}
\usepackage{tikz}
\usepackage{booktabs}
\usepackage{array}
\usepackage{listings}
\usepackage{xcolor}


% Set graphics path for images (relative to slides directory)
\graphicspath{{../Images/}}

% Title information
\title[Chapter 2.10: Automata]{Structured Computer Architecture\\Chapter 2.10: Automata\\Second-Order Systems (2-OS)}
\subtitle{An Order-Based Approach}
\author{John Glossner and Gheorghe M. Ştefan}
\institute{Based on the textbook}
\date{\today}

% Custom commands
\newcommand{\concept}[1]{\textbf{\color{accentcolor}#1}}
\newcommand{\defterm}[1]{\textit{#1}}

\begin{document}

% Title slide
\begin{frame}
\titlepage
\end{frame}

% Table of contents
\begin{frame}{Outline}
\tableofcontents
\end{frame}

% Learning Objectives
\begin{frame}{Learning Objectives}
By the end of this chapter, you will be able to:
\begin{itemize}
\item Understand second-order systems (2-OS) and the second feedback loop
\item Analyze finite state machines and their behavior
\item Design control automata for digital systems
\item Distinguish between simple and complex circuits
\item Understand the historical development of automata theory
\item Apply automata concepts to processor design
\end{itemize}
\end{frame}

% Section 1: Introduction to 2-OS
\section{Introduction to Second-Order Systems}

\begin{frame}{Evolution from 1-OS to 2-OS}
\begin{itemize}
\item \concept{1-OS Limitation}: Partial autonomy only
\begin{itemize}
\item State maintenance outside input signal range
\item No autonomous behavior evolution
\item Reactive to external stimuli
\end{itemize}
\vspace{0.5cm}
\item \concept{2-OS Innovation}: Second feedback loop
\begin{itemize}
\item Increases system autonomy
\item Enables autonomous behavior
\item Output evolution without input variation
\item Self-directed state transitions
\end{itemize}
\vspace{0.3cm}
\item \concept{Key Achievement}: 
\begin{itemize}
\item Behavior autonomy obtained
\item System can evolve independently
\item Foundation for intelligent systems
\end{itemize}
\end{itemize}
\end{frame}

\begin{frame}{Main Components of 2-OS}
\begin{itemize}
\item \concept{Control Automaton}:
\begin{itemize}
\item Manages system behavior
\item Implements state machine logic
\item Controls operation sequences
\end{itemize}
\vspace{0.5cm}
\item \concept{Program Counter (PC)}:
\begin{itemize}
\item Tracks instruction sequence
\item Provides autonomous progression
\item Enables program execution
\end{itemize}
\vspace{0.5cm}
\item \concept{Registers with Arithmetic \& Logic Unit (RALU)}:
\begin{itemize}
\item Data processing capability
\item Computational operations
\item State transformation logic
\end{itemize}
\end{itemize}
\end{frame}

% Section 2: Historical Development
\section{Historical Development of Automata}

\begin{frame}{Ancient Origins}
\begin{itemize}
\item \concept{Mythological Beginnings}:
\begin{itemize}
\item \textbf{Talos}: Bronze automaton from Greek mythology
\item Self-operating machines in ancient stories
\item Concept of artificial beings
\end{itemize}
\vspace{0.5cm}
\item \concept{Early Engineering}:
\begin{itemize}
\item \textbf{Hero of Alexandria (10-70 AD)}: Pneumatic and hydraulic devices
\item \textbf{Su Song (1020-1101 AD)}: Water-powered astronomical clock
\item Mechanical figurines and automated systems
\end{itemize}
\vspace{0.3cm}
\item \concept{Renaissance Innovations}:
\begin{itemize}
\item \textbf{Leonardo da Vinci (1452-1519)}: Mechanical knights
\item \textbf{Al-Jazari (1136-1206 AD)}: Sophisticated water clocks
\item Programmable features in mechanical systems
\end{itemize}
\end{itemize}
\end{frame}

\begin{frame}{Mathematical Foundations}
\begin{itemize}
\item \concept{Computational Pioneers}:
\begin{itemize}
\item \textbf{Pascal (17th century)}: Mechanical calculator
\item \textbf{Jacquard Loom (1804)}: Pioneered programmability
\item \textbf{Babbage \& Lovelace}: Analytical Engine concept
\end{itemize}
\vspace{0.5cm}
\item \concept{Theoretical Breakthroughs}:
\begin{itemize}
\item \textbf{Turing (1936)}: Turing Machine model
\item \textbf{von Neumann}: Self-replicating automata
\item \textbf{Chomsky}: Formal language classification
\end{itemize}
\vspace{0.3cm}
\item \concept{Modern Applications}:
\begin{itemize}
\item \textbf{Shannon}: Boolean circuits foundation
\item \textbf{Conway}: Game of Life cellular automata
\item AI, robotics, and formal verification
\end{itemize}
\end{itemize}
\end{frame}

% Section 3: Automata Definitions
\section{Automata Definitions and Classifications}
\begin{frame}{Generic Automaton Definition}
\begin{columns}
\begin{column}{0.6\textwidth}
\concept{Formal Definition}: Mathematical model of computation
\begin{itemize}
\item Set of states
\item Input alphabet
\item Transition function
\item Initial state
\item Accepting states
\end{itemize}
\vspace{0.3cm}
\concept{Two Main Categories}:
\begin{itemize}
\item \textbf{Small \& Complex}: Finite automata with intricate logic
\item \textbf{Large \& Simple}: Functional automata with regular patterns
\end{itemize}
\vspace{0.3cm}
\concept{Computational Power}:
\begin{itemize}
\item Language recognition capability
\item Problem-solving capacity
\item Complexity classification
\end{itemize}
\end{column}
\begin{column}{0.4\textwidth}
\begin{center}

\end{center}
\end{column}
\end{columns}
\end{frame}

\begin{frame}{Circuit Complexity Measures}
\begin{itemize}
\item \concept{Gate Size ($GS$)}:
\begin{itemize}
\item Total number of CMOS transistor pairs
\item Estimate of silicon area required
\item Hardware implementation cost
\end{itemize}
\vspace{0.3cm}
\item \concept{Area Size ($AS$)}:
\begin{itemize}
\item Actual silicon area dimension
\item Production cost factor
\item Yield considerations
\end{itemize}
\vspace{0.3cm}
\item \concept{Algorithmic Complexity ($C$)}:
\begin{itemize}
\item Minimal symbols needed for definition
\item Description length in HDL
\item Behavioral complexity measure
\end{itemize}
\end{itemize}
\end{frame}

\begin{frame}{Simple vs. Complex Circuits}
\begin{columns}
\begin{column}{0.5\textwidth}
\concept{Simple Circuits}:
\begin{itemize}
\item $C_{simple} << S_{simple}$
\item Complexity much smaller than size
\item Usually $C_{simple} \in O(1)$
\item Regular, repetitive patterns
\end{itemize}

\vspace{0.3cm}
\concept{Examples}:
\begin{itemize}
\item Adder arrays
\item Memory arrays
\item Prefix computation circuits
\item Regular interconnection patterns
\end{itemize}
\end{column}
\begin{column}{0.5\textwidth}
\concept{Complex Circuits}:
\begin{itemize}
\item $C_{complex} \sim S_{complex}$
\item Complexity similar to size
\item Irregular, random-like patterns
\item Difficult to describe concisely
\end{itemize}

\vspace{0.3cm}
\concept{Examples}:
\begin{itemize}
\item Random logic circuits
\item Irregular control logic
\item Custom optimization circuits
\item Application-specific designs
\end{itemize}
\end{column}
\end{columns}
\end{frame}

% Section 4: Finite State Machines
\section{Finite State Machines}

\begin{frame}{Finite Automaton Fundamentals}
\begin{itemize}
\item \concept{Formal Definition}: $FA = (Q, \Sigma, \delta, q_0, F)$
\begin{itemize}
\item $Q$: Finite set of states
\item $\Sigma$: Input alphabet
\item $\delta$: Transition function $Q \times \Sigma \rightarrow Q$
\item $q_0$: Initial state
\item $F$: Set of accepting states
\end{itemize}
\vspace{0.5cm}
\item \concept{Operation Principle}:
\begin{itemize}
\item Start in initial state $q_0$
\item Read input symbols sequentially
\item Transition according to $\delta$ function
\item Accept if final state is in $F$
\end{itemize}
\vspace{0.3cm}
\item \concept{Language Recognition}:
\begin{itemize}
\item Regular languages (Type 3)
\item Pattern matching capability
\item Lexical analysis applications
\end{itemize}
\end{itemize}
\end{frame}
\begin{frame}{State Transition Diagrams}
\begin{columns}
\begin{column}{0.6\textwidth}
\concept{Graphical Representation}:
\begin{itemize}
\item States as circles/nodes
\item Transitions as labeled arrows
\item Initial state marked with arrow
\item Accepting states double-circled
\end{itemize}
\vspace{0.3cm}
\concept{Design Process}:
\begin{enumerate}
\item Identify required states
\item Define input conditions
\item Specify state transitions
\item Mark accepting states
\item Verify completeness
\end{enumerate}
\vspace{0.3cm}
\concept{Analysis Techniques}:
\begin{itemize}
\item State reachability
\item Equivalence checking
\item Minimization algorithms
\item Determinism verification
\end{itemize}
\end{column}
\begin{column}{0.4\textwidth}
\begin{center}

\end{center}
\end{column}
\end{columns}
\end{frame}

\begin{frame}{Moore vs. Mealy Machines}
\begin{columns}
\begin{column}{0.5\textwidth}
\concept{Moore Machine}:
\begin{itemize}
\item Outputs depend only on current state
\item $\lambda: Q \rightarrow \Gamma$
\item Output changes with state transitions
\item Synchronous output behavior
\end{itemize}

\vspace{0.3cm}
\concept{Advantages}:
\begin{itemize}
\item Simpler timing analysis
\item Glitch-free outputs
\item Easier to design
\item More predictable behavior
\end{itemize}
\end{column}
\begin{column}{0.5\textwidth}
\concept{Mealy Machine}:
\begin{itemize}
\item Outputs depend on state and input
\item $\lambda: Q \times \Sigma \rightarrow \Gamma$
\item Faster response to inputs
\item More compact representation
\end{itemize}

\vspace{0.3cm}
\concept{Advantages}:
\begin{itemize}
\item Fewer states required
\item Immediate output response
\item More efficient implementation
\item Better for reactive systems
\end{itemize}
\end{column}
\end{columns}
\end{frame}

% Section 5: Control Automata
\section{Control Automata in Digital Systems}

\begin{frame}{Control Unit Design}
\begin{itemize}
\item \concept{Purpose}: Coordinate system operations
\begin{itemize}
\item Sequence control signals
\item Manage data flow
\item Implement instruction execution
\item Handle exceptions and interrupts
\end{itemize}
\vspace{0.5cm}
\item \concept{Design Approaches}:
\begin{itemize}
\item \textbf{Hardwired Control}: Fixed logic implementation
\item \textbf{Microprogrammed Control}: Stored program approach
\item \textbf{Hybrid Control}: Combination of both methods
\end{itemize}
\vspace{0.3cm}
\item \concept{Implementation Considerations}:
\begin{itemize}
\item Speed vs. flexibility trade-offs
\item Area and power constraints
\item Testability and verification
\item Modification and upgrade capability
\end{itemize}
\end{itemize}
\end{frame}

\begin{frame}{Program Counter (PC) Automaton}
\begin{itemize}
\item \concept{Function}: Autonomous instruction sequencing
\begin{itemize}
\item Maintains current instruction address
\item Provides sequential progression
\item Handles branches and jumps
\item Enables program execution flow
\end{itemize}
\vspace{0.5cm}
\item \concept{Operations}:
\begin{itemize}
\item \textbf{Increment}: Normal sequential execution
\item \textbf{Load}: Branch and jump operations
\item \textbf{Reset}: System initialization
\item \textbf{Hold}: Wait states and stalls
\end{itemize}
\vspace{0.3cm}
\item \concept{Advanced Features}:
\begin{itemize}
\item Branch prediction support
\item Exception handling
\item Multi-threading support
\item Debug and trace capabilities
\end{itemize}
\end{itemize}
\end{frame}

% Section 6: Counter-Expanded Automata
\section{Counter-Expanded Automata}

\begin{frame}{Beyond Regular Languages}
\begin{itemize}
\item \concept{Regular Language Limitation}:
\begin{itemize}
\item Cannot recognize $\{1^n0^n | n > 0\}$
\item Finite memory constraint
\item Need for counting capability
\end{itemize}
\vspace{0.5cm}
\item \concept{Counter-Expanded Automaton (CEA)}:
\begin{itemize}
\item Finite automaton + reversible counter
\item Counter provides zero flag to automaton
\item Enables context-free language recognition
\item 3-OS system (FA + Counter in loop)
\end{itemize}
\vspace{0.3cm}
\item \concept{Enhanced Capabilities}:
\begin{itemize}
\item Balanced parentheses checking
\item Nested structure recognition
\item Stack-like behavior simulation
\item Context-free grammar parsing
\end{itemize}
\end{itemize}
\end{frame}

\begin{frame}{CEA Example: $a^nb^n$ Language}
\begin{itemize}
\item \concept{Problem}: Recognize strings with equal numbers of $a$s and $b$s
\begin{itemize}
\item Context-free language
\item Requires counting mechanism
\item Cannot be solved by finite automaton alone
\end{itemize}
\vspace{0.5cm}
\item \concept{CEA Solution}:
\begin{itemize}
\item States: $\{q_0, q_1, q_2, q_3, q_4\}$
\item Counter operations: $\{nop, up, down\}$
\item Zero flag from counter
\item Delimiter symbol $e$ for well-formed strings
\end{itemize}
\vspace{0.3cm}
\item \concept{Operation Sequence}:
\begin{enumerate}
\item Count $a$s by incrementing counter
\item Switch to $b$ processing mode
\item Count $b$s by decrementing counter
\item Accept if counter reaches zero
\end{enumerate}
\end{itemize}
\end{frame}

% Section 7: Implementation Aspects
\section{Implementation and Design Considerations}

\begin{frame}{Hardware Implementation}
\begin{itemize}
\item \concept{State Encoding}:
\begin{itemize}
\item Binary encoding: Minimal flip-flops
\item One-hot encoding: Fast decoding
\item Gray code: Reduced switching
\item Custom encoding: Application-specific
\end{itemize}
\vspace{0.5cm}
\item \concept{Logic Implementation}:
\begin{itemize}
\item Combinational next-state logic
\item State register (flip-flops)
\item Output logic (Moore/Mealy)
\item Clock and reset distribution
\end{itemize}
\vspace{0.3cm}
\item \concept{Optimization Techniques}:
\begin{itemize}
\item State minimization
\item Logic synthesis
\item Timing optimization
\item Power reduction
\end{itemize}
\end{itemize}
\end{frame}

\begin{frame}{Design Verification}
\begin{itemize}
\item \concept{Functional Verification}:
\begin{itemize}
\item State coverage analysis
\item Transition coverage testing
\item Input sequence validation
\item Output correctness checking
\end{itemize}
\vspace{0.5cm}
\item \concept{Formal Methods}:
\begin{itemize}
\item Model checking techniques
\item Temporal logic verification
\item Equivalence checking
\item Property verification
\end{itemize}
\vspace{0.3cm}
\item \concept{Testing Strategies}:
\begin{itemize}
\item Directed test cases
\item Random test generation
\item Corner case analysis
\item Stress testing
\end{itemize}
\end{itemize}
\end{frame}

% Section 8: Applications
\section{Applications in Computer Architecture}

\begin{frame}{Processor Control Units}
\begin{itemize}
\item \concept{Instruction Decode}:
\begin{itemize}
\item Opcode interpretation
\item Addressing mode handling
\item Resource allocation
\item Pipeline control
\end{itemize}
\vspace{0.5cm}
\item \concept{Execution Control}:
\begin{itemize}
\item ALU operation sequencing
\item Memory access coordination
\item Register file management
\item Exception handling
\end{itemize}
\vspace{0.3cm}
\item \concept{System Control}:
\begin{itemize}
\item Interrupt processing
\item Power management
\item Debug support
\item Performance monitoring
\end{itemize}
\end{itemize}
\end{frame}

\begin{frame}{Memory Controllers}
\begin{itemize}
\item \concept{DRAM Control}:
\begin{itemize}
\item Refresh cycle management
\item Row/column addressing
\item Timing constraint enforcement
\item Bank management
\end{itemize}
\vspace{0.5cm}
\item \concept{Cache Controllers}:
\begin{itemize}
\item Hit/miss detection
\item Replacement policy implementation
\item Coherency protocol management
\item Prefetch control
\end{itemize}
\vspace{0.3cm}
\item \concept{I/O Controllers}:
\begin{itemize}
\item Protocol state machines
\item Data transfer management
\item Error detection and correction
\item Flow control
\end{itemize}
\end{itemize}
\end{frame}

% Section 9: Advanced Topics
\section{Advanced Automata Concepts}

\begin{frame}{Hierarchical State Machines}
\begin{itemize}
\item \concept{Motivation}: Managing complexity
\begin{itemize}
\item Large state spaces
\item Nested behaviors
\item Modular design
\item Reusability
\end{itemize}
\vspace{0.5cm}
\item \concept{Statecharts}:
\begin{itemize}
\item Hierarchical states
\item Concurrent regions
\item History states
\item Broadcast communication
\end{itemize}
\vspace{0.3cm}
\item \concept{Benefits}:
\begin{itemize}
\item Reduced complexity
\item Better maintainability
\item Improved readability
\item Easier verification
\end{itemize}
\end{itemize}
\end{frame}

\begin{frame}{Concurrent and Parallel Automata}
\begin{itemize}
\item \concept{Parallel Composition}:
\begin{itemize}
\item Multiple automata running simultaneously
\item Independent state evolution
\item Synchronization points
\item Communication mechanisms
\end{itemize}
\vspace{0.5cm}
\item \concept{Synchronization Methods}:
\begin{itemize}
\item Shared signals
\item Message passing
\item Rendezvous points
\item Barrier synchronization
\end{itemize}
\vspace{0.3cm}
\item \concept{Applications}:
\begin{itemize}
\item Multi-core processors
\item Distributed systems
\item Protocol implementations
\item Parallel algorithms
\end{itemize}
\end{itemize}
\end{frame}

% Section 10: Summary
\section{Chapter Summary}

\begin{frame}{Key Concepts Review}
\begin{itemize}
\item \concept{Second-Order Systems (2-OS)}:
\begin{itemize}
\item Second feedback loop enables autonomy
\item Behavior evolution without input changes
\item Foundation for intelligent systems
\end{itemize}
\vspace{0.3cm}
\item \concept{Automata Theory}:
\begin{itemize}
\item Mathematical models of computation
\item State-based behavior description
\item Language recognition capabilities
\end{itemize}
\vspace{0.3cm}
\item \concept{Implementation Aspects}:
\begin{itemize}
\item Hardware realization techniques
\item Optimization strategies
\item Verification methodologies
\end{itemize}
\end{itemize}
\end{frame}

\begin{frame}{Design Principles}
\begin{itemize}
\item \concept{Complexity Management}:
\begin{itemize}
\item Simple vs. complex circuit classification
\item Hierarchical design approaches
\item Modular implementation strategies
\end{itemize}
\vspace{0.3cm}
\item \concept{Performance Considerations}:
\begin{itemize}
\item State encoding optimization
\item Timing constraint satisfaction
\item Power and area trade-offs
\end{itemize}
\vspace{0.3cm}
\item \concept{Verification and Testing}:
\begin{itemize}
\item Formal verification methods
\item Coverage-driven testing
\item Model checking techniques
\end{itemize}
\end{itemize}
\end{frame}

\begin{frame}{Connection to Higher-Order Systems}
\begin{itemize}
\item \concept{Building Blocks for 3-OS (Processors)}:
\begin{itemize}
\item Control unit implementation
\item Instruction sequencing logic
\item Exception handling mechanisms
\end{itemize}
\vspace{0.3cm}
\item \concept{System-Level Integration}:
\begin{itemize}
\item Memory controller design
\item I/O protocol implementation
\item Inter-component communication
\end{itemize}
\vspace{0.3cm}
\item \concept{Advanced Applications}:
\begin{itemize}
\item Multi-core coordination
\item Distributed system control
\item Adaptive system behavior
\end{itemize}
\end{itemize}
\end{frame}

% Questions slide
\begin{frame}{Questions and Discussion}
\begin{center}
\Large
Questions?

\vspace{1cm}

\normalsize
\textbf{Discussion Topics:}
\begin{itemize}
\item How does the second feedback loop enable autonomous behavior?
\item What are the trade-offs between Moore and Mealy machines?
\item How do counter-expanded automata extend computational power?
\item What role do automata play in modern processor design?
\end{itemize}

\vspace{0.5cm}
\textbf{Next:} Chapter 2.11 - Processors: Third-Order Systems (3-OS)
\end{center}
\end{frame}

\end{document}

